% Appendix: Detailed Derivations
% Complete algebraic derivations referenced from main text

\appendix
\chapter{公式推导}

\section{差值均值统计量的方差分解}\label{sec:frt-algebra-derivation}

\textbf{背景}：在 \cref{eq:frt-algebra} 中，我们声称
\[
(n-1)s^2 = \sum_{Z_i=1} \{Y_i - \hat{\bar{Y}}(1)\}^2 + \sum_{Z_i=0} \{Y_i - \hat{\bar{Y}}(0)\}^2 + \frac{n_1 n_0}{n} \hat{\tau}^2 .
\tag{3.7}
\]
本节给出完整的代数推导。

\textbf{第一步}：从定义出发。
\[
(n-1)s^2 = \sumn (Y_i - \bar{Y})^2 = \sumn \big[ (Y_i - \bar{Y}) \big]^2 .
\]

\textbf{第二步}：对每个单元 $i$，将偏差 $Y_i - \bar{Y}$ 写成：
\[
Y_i - \bar{Y} = \underbrace{(Y_i - \hat{\bar{Y}}(Z_i))}_{\text{组内偏差}} + \underbrace{(\hat{\bar{Y}}(Z_i) - \bar{Y})}_{\text{组间偏差}} .
\]

\textbf{第三步}：平方展开：
\[
(Y_i - \bar{Y})^2 = (Y_i - \hat{\bar{Y}}(Z_i))^2 + (\hat{\bar{Y}}(Z_i) - \bar{Y})^2 + 2(Y_i - \hat{\bar{Y}}(Z_i))(\hat{\bar{Y}}(Z_i) - \bar{Y}) .
\]

\textbf{第四步}：对所有单元求和。关键观察是交叉项求和为零：
\[
\sum_{i=1}^n 2(Y_i - \hat{\bar{Y}}(Z_i))(\hat{\bar{Y}}(Z_i) - \bar{Y}) = 2 \sum_{i=1}^n (Y_i - \hat{\bar{Y}}(Z_i)) \cdot (\hat{\bar{Y}}(Z_i) - \bar{Y}) .
\]

当 $Z_i = 1$ 时，这个和化为：
\[
2 \sum_{Z_i=1} (Y_i - \hat{\bar{Y}}(1)) \cdot (\hat{\bar{Y}}(1) - \bar{Y}) = 2 (\hat{\bar{Y}}(1) - \bar{Y}) \underbrace{\sum_{Z_i=1} (Y_i - \hat{\bar{Y}}(1))}_{=0} = 0 ,
\]
因为 $\sum_{Z_i=1} Y_i = n_1 \hat{\bar{Y}}(1)$，所以 $\sum_{Z_i=1} (Y_i - \hat{\bar{Y}}(1)) = 0$。

同理，当 $Z_i = 0$ 时：
\[
2 \sum_{Z_i=0} (Y_i - \hat{\bar{Y}}(0)) \cdot (\hat{\bar{Y}}(0) - \bar{Y}) = 2 (\hat{\bar{Y}}(0) - \bar{Y}) \underbrace{\sum_{Z_i=0} (Y_i - \hat{\bar{Y}}(0))}_{=0} = 0 .
\]

因此交叉项全部抵消。

\textbf{第五步}：整理得到：
\[
\sumn (Y_i - \bar{Y})^2 = \underbrace{\sum_{Z_i=1} (Y_i - \hat{\bar{Y}}(1))^2 + \sum_{Z_i=0} (Y_i - \hat{\bar{Y}}(0))^2}_{\text{组内平方和}} + \underbrace{\sum_{i=1}^n (\hat{\bar{Y}}(Z_i) - \bar{Y})^2}_{\text{组间平方和}} .
\]

\textbf{第六步}：计算组间平方和：
\[
\sum_{i=1}^n (\hat{\bar{Y}}(Z_i) - \bar{Y})^2 = n_1 (\hat{\bar{Y}}(1) - \bar{Y})^2 + n_0 (\hat{\bar{Y}}(0) - \bar{Y})^2 .
\]

利用以下两个事实：
\begin{enumerate}
  \item $\hat{\bar{Y}}(1) - \hat{\bar{Y}}(0) = \hat{\tau}$
  \item $n_1 \hat{\bar{Y}}(1) + n_0 \hat{\bar{Y}}(0) = n \bar{Y}$（加权平均的性质）
\end{enumerate}

展开并化简：
\[
\begin{aligned}
& n_1 (\hat{\bar{Y}}(1) - \bar{Y})^2 + n_0 (\hat{\bar{Y}}(0) - \bar{Y})^2 \\
&= n_1 \hat{\bar{Y}}(1)^2 - 2 n_1 \hat{\bar{Y}}(1)\bar{Y} + n_1 \bar{Y}^2 + n_0 \hat{\bar{Y}}(0)^2 - 2 n_0 \hat{\bar{Y}}(0)\bar{Y} + n_0 \bar{Y}^2 \\
&= n_1 \hat{\bar{Y}}(1)^2 + n_0 \hat{\bar{Y}}(0)^2 - n \bar{Y}^2 .
\end{aligned}
\]

利用 $(n_1 \hat{\bar{Y}}(1) + n_0 \hat{\bar{Y}}(0))^2 = n^2 \bar{Y}^2$，展开得：
\[
n_1^2 \hat{\bar{Y}}(1)^2 + 2 n_1 n_0 \hat{\bar{Y}}(1)\hat{\bar{Y}}(0) + n_0^2 \hat{\bar{Y}}(0)^2 = n^2 \bar{Y}^2 .
\]

代入上式并利用 $\hat{\tau} = \hat{\bar{Y}}(1) - \hat{\bar{Y}}(0)$，最终化简为：
\[
n_1 (\hat{\bar{Y}}(1) - \bar{Y})^2 + n_0 (\hat{\bar{Y}}(0) - \bar{Y})^2 = \frac{n_1 n_0}{n} \hat{\tau}^2 .
\]

\textbf{结论}：代入即得 \cref{eq:frt-algebra}。

\section{Wilcoxon 秩和统计量的方差推导}\label{sec:wilcoxon-variance-derivation}

\textbf{背景}：在 \ref{ex:wilcoxon-mean-var} 中，我们声称在 $\HF$ 下
\[
\text{var}(W) = \frac{n_1 n_0 (n+1)}{12} .
\]
本节给出完整的代数推导。

\textbf{第一步}：回忆 $W = \sumn Z_i R_i$，其中 $R_i$ 是 $Y_i$ 在合并样本中的秩（固定值），$Z_i$ 是随机化的治疗指示变量。

\textbf{方差分解}：
\[
\text{var}(W) = \text{var}\left( \sumn Z_i R_i \right) = \sumn R_i^2 \, \text{var}(Z_i) + \sum_{i \neq j} R_i R_j \, \text{cov}(Z_i, Z_j) .
\]

\textbf{计算 $\text{var}(Z_i)$}：在 CRE 中，$Z_i \sim \text{Bernoulli}(n_1/n)$，所以
\[
\text{var}(Z_i) = \frac{n_1}{n} \cdot \frac{n_0}{n} = \frac{n_1 n_0}{n^2} .
\]

\textbf{计算 $\text{cov}(Z_i, Z_j)$ for $i \neq j$}：由于 $\sumn Z_i = n_1$（固定），我们有
\[
\text{cov}(Z_i, Z_j) = -\frac{n_1 n_0}{n^2(n-1)} \quad (i \neq j) .
\]

\textbf{代入}：
\[
\begin{aligned}
\text{var}(W) &= \left( \sumn R_i^2 \right) \frac{n_1 n_0}{n^2} + \sum_{i \neq j} R_i R_j \left( -\frac{n_1 n_0}{n^2(n-1)} \right) \\
&= \frac{n_1 n_0}{n^2} \left[ \sumn R_i^2 - \frac{1}{n-1} \sum_{i \neq j} R_i R_j \right] .
\end{aligned}
\]

注意 $\sum_{i \neq j} R_i R_j = \left( \sumn R_i \right)^2 - \sumn R_i^2$，代入得：
\[
\text{var}(W) = \frac{n_1 n_0}{n^2} \left[ \sumn R_i^2 - \frac{1}{n-1} \left( \left( \sumn R_i \right)^2 - \sumn R_i^2 \right) \right] .
\]

\textbf{计算 $\sumn R_i$ 和 $\sumn R_i^2$}：假设无 ties，则 $\{R_i\}$ 是 $1, 2, \ldots, n$ 的一个排列。
\[
\sumn R_i = 1 + 2 + \cdots + n = \frac{n(n+1)}{2} .
\]
\[
\sumn R_i^2 = 1^2 + 2^2 + \cdots + n^2 = \frac{n(n+1)(2n+1)}{6} .
\]

代入并化简：
\[
\begin{aligned}
\text{var}(W) &= \frac{n_1 n_0}{n^2} \left[ \frac{n(n+1)(2n+1)}{6} - \frac{1}{n-1} \left( \frac{n^2(n+1)^2}{4} - \frac{n(n+1)(2n+1)}{6} \right) \right] \\
&= \frac{n_1 n_0 (n+1)}{12} .
\end{aligned}
\]

\textbf{结论}：得证。

\section{女士品茶实验的超几何分布}\label{sec:lady-tea-hypergeometric}

\textbf{背景}：在 \ref{ex:lady-tea-hypergeometric} 中，我们讨论了 Fisher 的女士品茶实验。本节给出超几何分布的完整推导。

\textbf{设置}：
\begin{itemize}
  \item 8 杯茶中有 4 杯是"先加奶"($K=4$)，4 杯是"先加茶"
  \item 女士随机猜测其中 4 杯为"先加奶"
  \item 记 $X$ 为她猜对"先加奶"的杯数
\end{itemize}

\textbf{为什么是超几何分布？}

在 $\HF$ 下，女士的猜测是固定的（她选 4 杯作为"先加奶"的猜测）。实验者随机排列 8 杯茶的顺序呈现给女士。这等价于：从含有 $K=4$ 个"成功"的总体 $N=8$ 中，不放回抽取 $n=4$ 个，$X$ 是抽到的成功个数。

\textbf{超几何分布的概率质量函数}：
\[
\mathbb{P}(X = k) = \frac{\binom{K}{k} \binom{N-K}{n-k}}{\binom{N}{n}}, \quad k = 0, 1, 2, 3, 4 .
\]

\textbf{验证}：超几何分布的概率之和为 1。
\[
\sum_{k=0}^{4} \mathbb{P}(X = k) = \frac{1}{\binom{8}{4}} \sum_{k=0}^{4} \binom{4}{k} \binom{4}{4-k} = \frac{1}{70} \cdot \binom{8}{4} = 1 ,
\]
其中利用了组合恒等式 $\sum_{k} \binom{a}{k} \binom{b}{n-k} = \binom{a+b}{n}$。

\textbf{各概率的数值}：
\[
\begin{aligned}
\mathbb{P}(X = 0) &= \frac{\binom{4}{0}\binom{4}{4}}{\binom{8}{4}} = \frac{1 \cdot 1}{70} = \frac{1}{70}, \\
\mathbb{P}(X = 1) &= \frac{\binom{4}{1}\binom{4}{3}}{\binom{8}{4}} = \frac{4 \cdot 4}{70} = \frac{16}{70}, \\
\mathbb{P}(X = 2) &= \frac{\binom{4}{2}\binom{4}{2}}{\binom{8}{4}} = \frac{6 \cdot 6}{70} = \frac{36}{70}, \\
\mathbb{P}(X = 3) &= \frac{\binom{4}{3}\binom{4}{1}}{\binom{8}{4}} = \frac{4 \cdot 4}{70} = \frac{16}{70}, \\
\mathbb{P}(X = 4) &= \frac{\binom{4}{4}\binom{4}{0}}{\binom{8}{4}} = \frac{1 \cdot 1}{70} = \frac{1}{70}.
\end{aligned}
\]

\textbf{与二项分布的比较}：如果"抽取放回"（即女士每次猜测后被告知答案），则 $X \sim \text{Binomial}(4, 1/2)$，此时
\[
\mathbb{P}(X = 4) = \binom{4}{4} (1/2)^4 = 1/16 = 0.0625 > 0.014 .
\]

超几何分布（不放回）的尾部概率更小，使得 Fisher 的检验在零假设下更精确——这正是 FRT"精确"的含义。

\section{差值均值统计量的期望与方差推导}\label{sec:tau-hat-variance-derivation}

\textbf{背景}：在 \cref{eq:teststatistic} 中，我们定义了差值均值统计量
\[
\hat{\tau} = \hat{\bar{Y}}(1) - \hat{\bar{Y}}(0) .
\]
本节推导在 $\HF$ 下 $\hat{\tau}$ 的期望和方差。

\subsection*{期望的推导}

回忆在 CRE 中，治疗分配 $\mathbf{Z}$ 与潜在结果 $\{\mathbf{Y}(1), \mathbf{Y}(0)\}$ 独立，且每个单元被分配到治疗组的概率为 $n_1/n$：
\[
\mathbb{E}(Z_i) = \frac{n_1}{n}, \quad \mathbb{E}(1 - Z_i) = \frac{n_0}{n} .
\]

因此
\begin{align*}
\mathbb{E}(\hat{\tau})
&= \mathbb{E}\left( n_1^{-1} \sum_{i=1}^n Z_i Y_i - n_0^{-1} \sum_{i=1}^n (1 - Z_i) Y_i \right) \\
&= n_1^{-1} \sum_{i=1}^n \mathbb{E}(Z_i) Y_i - n_0^{-1} \sum_{i=1}^n \mathbb{E}(1 - Z_i) Y_i \\
&= n_1^{-1} \sum_{i=1}^n \frac{n_1}{n} Y_i - n_0^{-1} \sum_{i=1}^n \frac{n_0}{n} Y_i \\
&= \frac{1}{n} \sum_{i=1}^n Y_i - \frac{1}{n} \sum_{i=1}^n Y_i = 0 .
\end{align*}

\subsection*{方差的推导}

\textbf{第一步}：简化表达式。注意到
\[
\hat{\tau} = n_1^{-1} \sum_{i=1}^n Z_i Y_i - n_0^{-1} \sum_{i=1}^n (1 - Z_i) Y_i = \frac{n}{n_0} \cdot n_1^{-1} \sum_{i=1}^n Z_i Y_i - n_0^{-1} \sum_{i=1}^n Y_i .
\]

由于 $\sum_{i=1}^n Y_i$ 是常数（不依赖 $\mathbf{Z}$），我们只需计算
\[
\text{var}(\hat{\tau}) = \frac{n^2}{n_0^2 n_1^2} \cdot \text{var}\left( \sum_{i=1}^n Z_i Y_i \right) .
\]

\textbf{第二步}：计算 $\text{var}(Z_i)$ 和 $\text{cov}(Z_i, Z_j)$。在 CRE 中，
\[
Z_i \sim \text{Bernoulli}\left(\frac{n_1}{n}\right), \quad \text{var}(Z_i) = \frac{n_1}{n} \cdot \frac{n_0}{n} = \frac{n_1 n_0}{n^2} .
\]

由于 $\sum_{i=1}^n Z_i = n_1$（固定），对于 $i \neq j$，
\[
\text{cov}(Z_i, Z_j) = -\frac{n_1 n_0}{n^2 (n-1)} .
\]

\textbf{第三步}：计算 $\text{var}\left( \sum_{i=1}^n Z_i Y_i \right)$：
\begin{align*}
\text{var}\left( \sum_{i=1}^n Z_i Y_i \right)
&= \sum_{i=1}^n Y_i^2 \, \text{var}(Z_i) + \sum_{i \neq j} Y_i Y_j \, \text{cov}(Z_i, Z_j) \\
&= \left( \sum_{i=1}^n Y_i^2 \right) \frac{n_1 n_0}{n^2} + \sum_{i \neq j} Y_i Y_j \left( -\frac{n_1 n_0}{n^2 (n-1)} \right) .
\end{align*}

注意 $\sum_{i \neq j} Y_i Y_j = \left( \sum_{i=1}^n Y_i \right)^2 - \sum_{i=1}^n Y_i^2 = n^2 \bar{Y}^2 - \sum_{i=1}^n Y_i^2$，代入得：
\begin{align*}
\text{var}\left( \sum_{i=1}^n Z_i Y_i \right)
&= \frac{n_1 n_0}{n^2} \sum_{i=1}^n Y_i^2 - \frac{n_1 n_0}{n^2 (n-1)} \left( n^2 \bar{Y}^2 - \sum_{i=1}^n Y_i^2 \right) \\
&= \frac{n_1 n_0}{n^2 (n-1)} \left[ (n-1) \sum_{i=1}^n Y_i^2 - n^2 \bar{Y}^2 + \sum_{i=1}^n Y_i^2 \right] \\
&= \frac{n_1 n_0}{n^2 (n-1)} \left[ n \sum_{i=1}^n Y_i^2 - n^2 \bar{Y}^2 \right] \\
&= \frac{n_1 n_0}{n (n-1)} \left( \sum_{i=1}^n Y_i^2 - n \bar{Y}^2 \right) .
\end{align*}

\textbf{第四步}：利用 $s^2$ 的定义：
\[
s^2 = (n-1)^{-1} \sum_{i=1}^n (Y_i - \bar{Y})^2 = (n-1)^{-1} \left( \sum_{i=1}^n Y_i^2 - n \bar{Y}^2 \right) .
\]

因此
\[
\text{var}\left( \sum_{i=1}^n Z_i Y_i \right) = \frac{n_1 n_0}{n} \cdot s^2 .
\]

\textbf{第五步}：代回 $\text{var}(\hat{\tau})$：
\[
\text{var}(\hat{\tau}) = \frac{n^2}{n_0^2 n_1^2} \cdot \frac{n_1 n_0}{n} \cdot s^2 = \frac{n}{n_1 n_0} \cdot s^2 .
\]

\textbf{结论}：
\[
\boxed{ \text{var}(\hat{\tau}) = \frac{n}{n_1 n_0} \cdot s^2 } .
\]

\textbf{直观理解}：
\begin{itemize}
  \item 方差与总样本量 $n$ 成正比——样本越大，估计越精确
  \item 方差与 $n_1 n_0$ 成反比——当 $n_1 = n_0 = n/2$ 时方差最小
  \item 这与直觉一致：治疗组和对照组越平衡，我们越能准确估计因果效应
\end{itemize}

\section{CRE 中分层平衡性的推导}\label{sec:appendix-5-2-balance}

\textbf{背景（Background）}：在完全随机实验（CRE）中，我们固定总的治疗组和对照组样本量 $n_1$ 和 $n_0$，但各层内的分配数量 $n_{[k]1}$ 和 $n_{[k]0}$ 是随机的。本节证明各层治疗组和对照组比例之差的期望为零。

\textbf{参数定义（Parameter Definitions）}：
\begin{itemize}
  \item $n$：总单元数
  \item $n_{[k]}$：第 $k$ 层的单元数
  \item $n_1$：总治疗组单元数
  \item $n_0$：总对照组单元数（$n_0 = n - n_1$）
  \item $n_{[k]1}$：第 $k$ 层中分配到治疗组的单元数
  \item $n_{[k]0}$：第 $k$ 层中分配到对照组的单元数（$n_{[k]0} = n_{[k]} - n_{[k]1}$）
  \item $\pi_{[k]} = n_{[k]}/n$：第 $k$ 层在总体中的比例
\end{itemize}

\textbf{已知条件（Given）}：
\begin{itemize}
  \item 在 CRE 下，治疗分配向量 $\mathbf{Z}$ 是从所有 $\binom{n}{n_1}$ 种分配中等概率随机选取的
  \item 各单元被分配到治疗组的概率为 $n_1/n$
\end{itemize}

\textbf{目标（Goal）}：证明
\[
\mathbb{E}\left( \frac{n_{[k]1}}{n_1} - \frac{n_{[k]0}}{n_0} \right) = 0.
\]

\textbf{推导步骤（Derivation Steps）}：

\textbf{第一步}：建立 $n_{[k]1}$ 的期望。

在 CRE 下，每个单元被分配到治疗组的概率为 $n_1/n$。因此，对于第 $k$ 层中的任意单元 $i$（满足 $X_i = k$）：
\[
\mathbb{E}(Z_i) = \frac{n_1}{n}.
\]

第 $k$ 层中分配到治疗组的单元数为：
\[
n_{[k]1} = \sum_{i: X_i = k} Z_i.
\]

由于该层有 $n_{[k]}$ 个单元，且 $Z_i$ 之间是负相关的（总和固定为 $n_1$），我们有：
\[
\mathbb{E}(n_{[k]1}) = \sum_{i: X_i = k} \mathbb{E}(Z_i) = n_{[k]} \cdot \frac{n_1}{n} = \pi_{[k]} \cdot n_1.
\]

\textbf{第二步}：计算比例的期望。

由于 $n_{[k]0} = n_{[k]} - n_{[k]1}$，我们有：
\[
\mathbb{E}(n_{[k]0}) = n_{[k]} - \mathbb{E}(n_{[k]1}) = n_{[k]} - \pi_{[k]} n_1 = n_{[k]} - \frac{n_{[k]}}{n} n_1 = n_{[k]} \left(1 - \frac{n_1}{n}\right) = n_{[k]} \cdot \frac{n_0}{n} = \pi_{[k]} \cdot n_0.
\]

\textbf{第三步}：计算比例之差的期望。

\[
\mathbb{E}\left( \frac{n_{[k]1}}{n_1} - \frac{n_{[k]0}}{n_0} \right) = \frac{\mathbb{E}(n_{[k]1})}{n_1} - \frac{\mathbb{E}(n_{[k]0})}{n_0} = \frac{\pi_{[k]} n_1}{n_1} - \frac{\pi_{[k]} n_0}{n_0} = \pi_{[k]} - \pi_{[k]} = 0.
\]

\textbf{结论}：得证 $\mathbb{E}\left( n_{[k]1}/n_1 - n_{[k]0}/n_0 \right) = 0$。

\textbf{直观理解}：这个结果表明，虽然在一次实现中 $n_{[k]1}/n_1 - n_{[k]0}/n_0$ 可能显著非零（由于随机波动），但这种波动在期望意义下是平衡的。随机化保证了渐近公平性——各层在治疗组和对照组中的比例在期望上与总体比例一致。

\section{分层估计量的期望与方差推导}\label{sec:appendix-5-8}

\textbf{背景（Background）}：在分层随机实验（SRE）中，我们使用学生化分层估计量
\[
t_\SRE = \frac{\hat{\tau}_\SRE}{\sqrt{\hat{V}_\SRE}}
\]
作为 Fisher 随机化检验的统计量。本节推导在 sharp null $H_{0\Fisher}: Y_i(1) = Y_i(0)$ 下 $t_\SRE$ 的期望和方差。

\textbf{参数定义（Parameter Definitions）}：
\begin{itemize}
  \item $\hat{\tau}_\SRE = \sum_{k=1}^K \pi_{[k]} \hat{\tau}_{[k]}$：分层估计量
  \item $\hat{\tau}_{[k]} = \bar{Y}_{[k]}(1) - \bar{Y}_{[k]}(0)$：第 $k$ 层的差值均值
  \item $\hat{V}_\SRE = \sum_{k=1}^K \pi_{[k]}^2 \left( \frac{\hat{S}_{[k]}^2(1)}{n_{[k]1}} + \frac{\hat{S}_{[k]}^2(0)}{n_{[k]0}} \right)$：方差估计量
  \item $\hat{S}_{[k]}^2(1)$、$\hat{S}_{[k]}^2(0)$：第 $k$ 层内治疗组和对照组的样本方差
\end{itemize}

\textbf{目标（Goal）}：推导 $\hat{\tau}_\SRE$ 和 $\hat{V}_\SRE$ 在 $H_{0\Fisher}$ 下的期望和方差。

\textbf{推导步骤（Derivation Steps）}：

\textbf{第一步}：$\hat{\tau}_{[k]}$ 的期望。

在 $H_{0\Fisher}$ 下，所有 $Y_i(1) = Y_i(0) = Y_i$。因此，
\[
\hat{\tau}_{[k]} = \bar{Y}_{[k]}(1) - \bar{Y}_{[k]}(0) = \frac{\sum_{i: X_i=k} Z_i Y_i}{n_{[k]1}} - \frac{\sum_{i: X_i=k} (1-Z_i) Y_i}{n_{[k]0}}.
\]

利用 $\mathbb{E}(Z_i) = n_1/n$ 和 $\mathbb{E}(1-Z_i) = n_0/n$，以及 $Y_i$ 是固定值（不依赖 $Z$），我们有：
\[
\mathbb{E}(\hat{\tau}_{[k]}) = \frac{n_{[k]}}{n_{[k]1}} \cdot \frac{n_1}{n} \cdot \bar{Y}_{[k]} - \frac{n_{[k]}}{n_{[k]0}} \cdot \frac{n_0}{n} \cdot \bar{Y}_{[k]} = \bar{Y}_{[k]} \left( \frac{n_{[k]}}{n_{[k]1}} \cdot \frac{n_1}{n} - \frac{n_{[k]}}{n_{[k]0}} \cdot \frac{n_0}{n} \right).
\]

由于 $n_{[k]1}/n_{[k]} = e_{[k]}$（层别倾向得分）和 $n_{[k]0}/n_{[k]} = 1 - e_{[k]}$，以及 $n_1/n$ 和 $n_0/n$ 是常数，上述表达式可简化为 $\bar{Y}_{[k]} \cdot (\text{常数} - \text{常数}) = 0$。

更直接地：由于 $\mathbb{E}(Z_i) = n_1/n$ 对所有单元成立，在 $H_{0\Fisher}$ 下，
\[
\mathbb{E}(\hat{\tau}_{[k]}) = \frac{1}{n_{[k]1}} \sum_{i: X_i=k} \mathbb{E}(Z_i) Y_i - \frac{1}{n_{[k]0}} \sum_{i: X_i=k} \mathbb{E}(1-Z_i) Y_i = \frac{n_1}{n} \bar{Y}_{[k]} - \frac{n_0}{n} \bar{Y}_{[k]} = 0.
\]

\textbf{第二步}：$\hat{\tau}_\SRE$ 的期望。

由于 $\hat{\tau}_\SRE = \sum_{k=1}^K \pi_{[k]} \hat{\tau}_{[k]}$ 且 $\pi_{[k]}$ 是常数，
\[
\mathbb{E}(\hat{\tau}_\SRE) = \sum_{k=1}^K \pi_{[k]} \mathbb{E}(\hat{\tau}_{[k]}) = \sum_{k=1}^K \pi_{[k]} \cdot 0 = 0.
\]

\textbf{第三步}：$\hat{\tau}_{[k]}$ 的方差。

利用 Neyman (1923) 的方差公式（见 \cref{sec:tau-hat-variance-derivation}），在第 $k$ 层内，
\[
\text{var}(\hat{\tau}_{[k]}) = \frac{S_{[k]}^2(1)}{n_{[k]1}} + \frac{S_{[k]}^2(0)}{n_{[k]0}} - \frac{S_{[k]}^2(\tau)}{n_{[k]}}.
\]

在 $H_{0\Fisher}$ 下，$S_{[k]}^2(\tau) = 0$，所以
\[
\text{var}(\hat{\tau}_{[k]}) = \frac{S_{[k]}^2(1)}{n_{[k]1}} + \frac{S_{[k]}^2(0)}{n_{[k]0}}.
\]

\textbf{第四步}：$\hat{\tau}_\SRE$ 的方差。

由于各层之间的随机化是独立的，
\[
\text{var}(\hat{\tau}_\SRE) = \sum_{k=1}^K \pi_{[k]}^2 \text{var}(\hat{\tau}_{[k]}) = \sum_{k=1}^K \pi_{[k]}^2 \left( \frac{S_{[k]}^2(1)}{n_{[k]1}} + \frac{S_{[k]}^2(0)}{n_{[k]0}} \right).
\]

\textbf{第五步}：$\hat{V}_\SRE$ 的期望。

利用 $\mathbb{E}(\hat{S}_{[k]}^2(1)) = S_{[k]}^2(1)$ 和 $\mathbb{E}(\hat{S}_{[k]}^2(0)) = S_{[k]}^2(0)$，
\[
\mathbb{E}(\hat{V}_\SRE) = \sum_{k=1}^K \pi_{[k]}^2 \left( \frac{S_{[k]}^2(1)}{n_{[k]1}} + \frac{S_{[k]}^2(0)}{n_{[k]0}} \right) = \text{var}(\hat{\tau}_\SRE) + \sum_{k=1}^K \pi_{[k]}^2 \frac{S_{[k]}^2(\tau)}{n_{[k]}}.
\]

在 $H_{0\Fisher}$ 下，$S_{[k]}^2(\tau) = 0$，所以 $\mathbb{E}(\hat{V}_\SRE) = \text{var}(\hat{\tau}_\SRE)$。

\textbf{结论}：在 $H_{0\Fisher}$ 下，$\hat{\tau}_\SRE$ 的期望为零，方差为 $\hat{V}_\SRE$ 的期望。因此 $t_\SRE = \hat{\tau}_\SRE / \sqrt{\hat{V}_\SRE}$ 是一个合理的学生化统计量。

\section{分层 Neymanian 方差的推导}\label{sec:appendix-5-10}

\textbf{背景（Background）}：在分层随机实验（SRE）中，本节推导第 $k$ 层差值均值估计量 $\hat{\tau}_{[k]}$ 的方差公式，以及分层估计量 $\hat{\tau}_\SRE$ 的方差公式。

\textbf{参数定义（Parameter Definitions）}：
\begin{itemize}
  \item $S_{[k]}^2(1)$、$S_{[k]}^2(0)$：第 $k$ 层内潜在结果 $Y(1)$ 和 $Y(0)$ 的有限总体方差
  \item $S_{[k]}^2(\tau)$：第 $k$ 层内个体因果效应 $\tau_i = Y_i(1) - Y_i(0)$ 的有限总体方差
  \item $\hat{\tau}_{[k]} = \bar{Y}_{[k]}(1) - \bar{Y}_{[k]}(0)$：第 $k$ 层的差值均值
  \item $\hat{\tau}_\SRE = \sum_{k=1}^K \pi_{[k]} \hat{\tau}_{[k]}$：分层估计量
\end{itemize}

\textbf{目标（Goal）}：
\begin{enumerate}
  \item 证明 \eqref{eq:5-10}：$\text{var}(\hat{\tau}_{[k]}) = \frac{S_{[k]}^2(1)}{n_{[k]1}} + \frac{S_{[k]}^2(0)}{n_{[k]0}} - \frac{S_{[k]}^2(\tau)}{n_{[k]}}$
  \item 证明 \eqref{eq:5-11}：$\text{var}(\hat{\tau}_\SRE) = \sum_{k=1}^K \pi_{[k]}^2 \text{var}(\hat{\tau}_{[k]})$
\end{enumerate}

\textbf{推导步骤（Derivation Steps）}：

\subsection*{式 \eqref{eq:5-10} 的推导}

\textbf{第一步}：应用 Neyman (1923) 定理于第 $k$ 层。

第 $k$ 层本质上是一个独立的 CRE，只是总体大小为 $n_{[k]}$ 而非 $n$。直接应用 Neyman 的方差公式（见 \cref{sec:tau-hat-variance-derivation}）：
\[
\text{var}(\hat{\tau}_{[k]}) = \frac{S_{[k]}^2(1)}{n_{[k]1}} + \frac{S_{[k]}^2(0)}{n_{[k]0}} - \frac{S_{[k]}^2(\tau)}{n_{[k]}}.
\]

\textbf{直观解释}：这个方差公式与 CRE 中的方差公式结构相同，只是将总体量替换为层别总体量。

\subsection*{式 \eqref{eq:5-11} 的推导}

\textbf{第二步}：利用层间独立性。

在 SRE 中，各层内的治疗分配是独立进行的。因此，$\hat{\tau}_{[1]}, \ldots, \hat{\tau}_{[K]}$ 之间相互独立。

\textbf{第三步}：计算分层估计量的方差。

\[
\text{var}(\hat{\tau}_\SRE) = \text{var}\left( \sum_{k=1}^K \pi_{[k]} \hat{\tau}_{[k]} \right) = \sum_{k=1}^K \pi_{[k]}^2 \text{var}(\hat{\tau}_{[k]}) + \sum_{k \neq j} \pi_{[k]} \pi_{[j]} \text{cov}(\hat{\tau}_{[k]}, \hat{\tau}_{[j]}).
\]

由于各层之间的独立性，$\text{cov}(\hat{\tau}_{[k]}, \hat{\tau}_{[j]}) = 0$（对 $k \neq j$）。因此，
\[
\text{var}(\hat{\tau}_\SRE) = \sum_{k=1}^K \pi_{[k]}^2 \text{var}(\hat{\tau}_{[k]}).
\]

\textbf{结论}：代入 \eqref{eq:5-10} 即得
\[
\text{var}(\hat{\tau}_\SRE) = \sum_{k=1}^K \pi_{[k]}^2 \left( \frac{S_{[k]}^2(1)}{n_{[k]1}} + \frac{S_{[k]}^2(0)}{n_{[k]0}} - \frac{S_{[k]}^2(\tau)}{n_{[k]}} \right).
\]

\section{CRE 与 SRE 方差差异的推导}\label{sec:appendix-5-14}

\textbf{背景（Background）}：本节推导在常数倾向得分条件下，CRE 估计量 $\hat{\tau}$ 与 SRE 估计量 $\hat{\tau}_\SRE$ 方差差异的近似公式。

\textbf{参数定义（Parameter Definitions）}：
\begin{itemize}
  \item $\hat{\tau} = \bar{Y}(1) - \bar{Y}(0)$：CRE 的差值均值估计量
  \item $\hat{\tau}_\SRE = \sum_{k=1}^K \pi_{[k]} \hat{\tau}_{[k]}$：SRE 的分层估计量
  \item $S^2(1)$、$S^2(0)$：总体的潜在结果有限总体方差
  \item $S_{[k]}^2(1)$、$S_{[k]}^2(0)$：第 $k$ 层的潜在结果有限总体方差
  \item $\bar{Y}(1)$、$\bar{Y}(0)$：总体的潜在结果均值
  \item $\bar{Y}_{[k]}(1)$、$\bar{Y}_{[k]}(0)$：第 $k$ 层的潜在结果均值
\end{itemize}

\textbf{已知条件（Given）}：
\begin{itemize}
  \item 在 CRE 下：$\text{var}_{CRE}(\hat{\tau}) = \frac{S^2(1)}{n_1} + \frac{S^2(0)}{n_0} - \frac{S^2(\tau)}{n}$
  \item 在 SRE 下：$\text{var}_{SRE}(\hat{\tau}_\SRE) = \sum_{k=1}^K \pi_{[k]}^2 \left( \frac{S_{[k]}^2(1)}{n_{[k]1}} + \frac{S_{[k]}^2(0)}{n_{[k]0}} - \frac{S_{[k]}^2(\tau)}{n_{[k]}} \right)$
  \item 方差分析分解：$S^2(1) = \sum_{k=1}^K \pi_{[k]} S_{[k]}^2(1) + \sum_{k=1}^K \pi_{[k]} \{\bar{Y}_{[k]}(1) - \bar{Y}(1)\}^2 - \text{层内变异}$
\end{itemize}

\textbf{目标（Goal）}：证明方差差异约为
\[
\sum_{k=1}^K \frac{\pi_{[k]}}{n} \left\{ \sqrt{\frac{n_0}{n_1}} \{\bar{Y}_{[k]}(1) - \bar{Y}(1)\} + \sqrt{\frac{n_1}{n_0}} \{\bar{Y}_{[k]}(0) - \bar{Y}(0)\} \right\}^2 \geq 0.
\]

\textbf{推导步骤（Derivation Steps）}：

\textbf{第一步}：方差分析分解。

对 $S^2(1)$ 进行层间分解：
\[
\begin{aligned}
S^2(1) & = (n-1)^{-1} \sum_{i=1}^n \{Y_i(1) - \bar{Y}(1)\}^2 \\
& = \sum_{k=1}^K \frac{n_{[k]}-1}{n-1} \cdot \frac{1}{n_{[k]}-1} \sum_{i: X_i=k} \{Y_i(1) - \bar{Y}_{[k]}(1)\}^2 \\
& \quad + \sum_{k=1}^K \frac{n_{[k]}}{n-1} \{\bar{Y}_{[k]}(1) - \bar{Y}(1)\}^2.
\end{aligned}
\]

当 $n_{[k]}$ 较大时，$(n_{[k]}-1)/(n-1) \approx n_{[k]}/n = \pi_{[k]}$，于是
\[
S^2(1) \approx \sum_{k=1}^K \pi_{[k]} S_{[k]}^2(1) + \sum_{k=1}^K \pi_{[k]} \{\bar{Y}_{[k]}(1) - \bar{Y}(1)\}^2.
\]

同理对 $S^2(0)$ 和 $S^2(\tau)$ 进行分解。

\textbf{第二步}：代入 CRE 方差公式。

\[
\text{var}_{CRE}(\hat{\tau}) = \frac{1}{n_1} \sum_{k=1}^K \pi_{[k]} S_{[k]}^2(1) + \frac{1}{n_0} \sum_{k=1}^K \pi_{[k]} S_{[k]}^2(0) - \frac{1}{n} \sum_{k=1}^K \pi_{[k]} S_{[k]}^2(\tau) + \text{层间变异项}.
\]

\textbf{第三步}：识别差异项。

层间变异项为：
\[
\frac{1}{n_1} \sum_{k=1}^K \pi_{[k]} \{\bar{Y}_{[k]}(1) - \bar{Y}(1)\}^2 + \frac{1}{n_0} \sum_{k=1}^K \pi_{[k]} \{\bar{Y}_{[k]}(0) - \bar{Y}(0)\}^2 - \frac{1}{n} \sum_{k=1}^K \pi_{[k]} (\tau_{[k]} - \tau)^2.
\]

\textbf{第四步}：近似简化。

当层样本量较大时，利用 $n_{[k]1} \approx \pi_{[k]} n_1$ 和 $n_{[k]0} \approx \pi_{[k]} n_0$，层间变异项近似为：
\[
\sum_{k=1}^K \frac{\pi_{[k]}}{n} \left\{ \sqrt{\frac{n_0}{n_1}} \{\bar{Y}_{[k]}(1) - \bar{Y}(1)\} + \sqrt{\frac{n_1}{n_0}} \{\bar{Y}_{[k]}(0) - \bar{Y}(0)\} \right\}^2.
\]

\textbf{结论}：由于每一项都是平方和，该表达式非负。当且仅当对所有 $k$ 都有
\[
\sqrt{\frac{n_0}{n_1}} \{\bar{Y}_{[k]}(1) - \bar{Y}(1)\} + \sqrt{\frac{n_1}{n_0}} \{\bar{Y}_{[k]}(0) - \bar{Y}(0)\} = 0
\]
时（即协变量不能预测潜在结果），差异为零。这证明了 SRE 在协变量能预测结果时比 CRE 更有效。

\section{CRE 条件分布的推导}\label{sec:appendix-5-5}

\textbf{背景（Background）}：本节证明在给定各层治疗分配数量 $\mathbf{n} = \{n_{[k]1}, n_{[k]0}\}_{k=1}^K$ 的条件下，CRE 的治疗分配向量 $\mathbf{Z}$ 的条件分布与 SRE 的分布相同。

\textbf{参数定义（Parameter Definitions）}：
\begin{itemize}
  \item $\mathbf{Z} = (Z_1, \ldots, Z_n)$：治疗分配向量
  \item $n_{[k]1}$：第 $k$ 层中分配到治疗组的单元数
  \item $n_{[k]0}$：第 $k$ 层中分配到对照组的单元数
  \item $\mathbf{n} = \{n_{[k]1}, n_{[k]0}\}_{k=1}^K$：所有层的分配数量向量
\end{itemize}

\textbf{已知条件（Given）}：
\begin{itemize}
  \item 在 CRE 下：$\mathbb{P}(\mathbf{Z} = \mathbf{z}) = \frac{1}{\binom{n}{n_1}}$ 对所有满足 $\sum Z_i = n_1$ 的分配 $\mathbf{z}$
  \item 在 SRE 下：第 $k$ 层的分配概率为 $\frac{1}{\binom{n_{[k]}}{n_{[k]1}}}$
\end{itemize}

\textbf{目标（Goal）}：证明
\[
\Pr_{CRE}(\mathbf{Z} = \mathbf{z} \mid \mathbf{n}) = \frac{1}{\prod_{k=1}^K \binom{n_{[k]}}{n_{[k]1}}}.
\]

\textbf{推导步骤（Derivation Steps）}：

\textbf{第一步}：计算联合概率。

在 CRE 下，治疗分配向量 $\mathbf{Z}$ 的概率为：
\[
\Pr_{CRE}(\mathbf{Z} = \mathbf{z}) = \frac{1}{\binom{n}{n_1}}.
\]

\textbf{第二步}：计算条件概率。

给定 $\mathbf{n}$ 条件下 $\mathbf{Z} = \mathbf{z}$ 的条件概率为：
\[
\Pr_{CRE}(\mathbf{Z} = \mathbf{z} \mid \mathbf{n}) = \frac{\Pr_{CRE}(\mathbf{Z} = \mathbf{z}, \mathbf{n})}{\Pr_{CRE}(\mathbf{n})}.
\]

事件 $\{\mathbf{Z} = \mathbf{z}, \mathbf{n}\}$ 意味着第 $k$ 层恰好有 $n_{[k]1}$ 个单元被分配到治疗组。由于各层的分配是独立的，这个联合事件的概率为：
\[
\Pr_{CRE}(\mathbf{Z} = \mathbf{z}, \mathbf{n}) = \prod_{k=1}^K \frac{1}{\binom{n_{[k]}}{n_{[k]1}}}.
\]

\textbf{第三步}：计算边际概率。

为了计算 $\Pr_{CRE}(\mathbf{n})$，我们对所有满足第 $k$ 层有 $n_{[k]1}$ 个治疗组单元的分配 $\mathbf{z}$ 求和：
\[
\Pr_{CRE}(\mathbf{n}) = \sum_{\mathbf{z}: n_{[k]1}(\mathbf{z}) = n_{[k]1}} \Pr_{CRE}(\mathbf{Z} = \mathbf{z}) = \sum_{\mathbf{z}: n_{[k]1}(\mathbf{z}) = n_{[k]1}} \frac{1}{\binom{n}{n_1}}.
\]

满足该条件的分配数量为 $\prod_{k=1}^K \binom{n_{[k]}}{n_{[k]1}}$，因此
\[
\Pr_{CRE}(\mathbf{n}) = \prod_{k=1}^K \binom{n_{[k]}}{n_{[k]1}} \cdot \frac{1}{\binom{n}{n_1}}.
\]

\textbf{第四步}：代入条件概率公式。

\[
\Pr_{CRE}(\mathbf{Z} = \mathbf{z} \mid \mathbf{n}) = \frac{\prod_{k=1}^K \frac{1}{\binom{n_{[k]}}{n_{[k]1}}}}{\prod_{k=1}^K \binom{n_{[k]}}{n_{[k]1}} \cdot \frac{1}{\binom{n}{n_1}}} = \frac{1}{\prod_{k=1}^K \binom{n_{[k]}}{n_{[k]1}}}.
\]

\textbf{结论}：得证 CRE 在给定各层分配数量条件下的条件分布与 SRE 的分布相同。这建立了分层与后分层的对偶性：分层在设计阶段使用协变量，后分层在分析阶段使用协变量。
